% ===== UNIT 02: Motivation (1.1) =====
% Status: REWRITTEN  (change to REWRITTEN / REVIEWED when done; see CHECKLIST.md)
% Rules: style rewrite only. Keep every \cite, \ref, \label, number, equation and table value as is. Change the wording, not the content.
\subsection{Motivation}
Vehicle re-identification is heavily used in intelligent
transportation systems (ITS), supporting applications including urban traffic
surveillance, law enforcement and smart city infrastructure~\cite{liu2016veri,liu2018provid}.
The task requires matching the
same physical vehicle across multiple spatially distributed cameras whose
fields of view do not overlap. The same vehicle can look very different: it may be observed from different angles, under
different illumination conditions, at varying resolutions, and after arbitrary
time intervals~\cite{liu2016veri}.

The VeRi-776 dataset~\cite{liu2016veri,liu2018provid}, introduced by Liu
\emph{et al.} in 2016 and expanded by the PROVID framework published
in \emph{IEEE Transactions on Multimedia} in 2018, has become the
benchmark for vehicle ReID evaluation. It contains over 50\,000 images of 776
vehicles captured by 20 cameras covering a $1.0\,\text{km}^{2}$ urban area over
24 hours, with each vehicle appearing in 2 to 18 cameras under varied
viewpoints, illuminations, and occlusions~\cite{liu2016veri,liu2018provid}. The
standard evaluation protocol employs 37\,781 training images of 576
identities, 11\,579 gallery images, and 1\,678 query images across 200 test
identities~\cite{liu2018provid}. Performance is measured by mean Average
Precision (mAP) and Rank-1 Cumulative Matching Characteristic (CMC).
